\originalchapter{研究生作业七：作用量与粒子运动}
题面据8.311作业7的全部5题. 原页眉误写作业6, 文件日期为2004年3月31日. 归档名\texttt{solution7.pdf}实际页眉和内容均为作业8解答, 因而当前未取得与这套题匹配的官方解. 本章解答全部为AI独立推导, 不冒称官方题7翻译.
\originalsection{从作用量导出Lorentz力}
\begin{exercise}
原题1. 非相对论电荷$q$的$L=m\dot{\rr}^2/2+q\dot{\rr}\cdot\vct A-qV$, 用Euler--Lagrange求运动方程.
\end{exercise}
\begin{solution}
正则动量$P_i=mv_i+qA_i$. 沿轨迹微分$\dd A_i/\dd t=\partial_tA_i+v_j\partial_jA_i$, 而$\partial L/\partial x_i=qv_j\partial_iA_j-q\partial_iV$. 因而
$m\dot v_i=q[-\partial_iV-\partial_tA_i+v_j(\partial_iA_j-\partial_jA_i)]$. 最后括号是$(\vct v\times\nabla\times\vct A)_i$, 所以$m\dot{\vct v}=q(\ee+\vct v\times\bb)$. 若只对$A_i$取偏时间导数会漏掉磁力的一半结构.
\end{solution}
\originalsection{规范变换与作用量端点}
\begin{exercise}
原题2. $\vct A'=\vct A+\nabla\Lambda$, $V'=V-\partial_t\Lambda$. 求$L,S$变化并证明固定时空端点下方程不变.
\end{exercise}
\begin{solution}
$L'-L=q(\vct v\cdot\nabla\Lambda+\partial_t\Lambda)=q\dd\Lambda(\rr(t),t)/\dd t$. 因而$S'-S=q[\Lambda(\rr_2,t_2)-\Lambda(\rr_1,t_1)]$. 作用量数值一般改变, 但变化只依赖固定端点, 其路径变分零, Euler--Lagrange方程不变. 这解释为何同一物理场的不同势可描述相同轨迹.
\end{solution}
\originalsection{正则形式的独立变分}
\begin{exercise}
原题3. $S=\int[\sum p_i\dot q_i-H(q,p,t)]\dd t$, $q,p$独立变分且仅$\delta q$端点为零. 导出正则方程.
\end{exercise}
\begin{solution}
对$p_i\delta\dot q_i$分部积分, 边界项$p_i\delta q_i|_{t_1}^{t_2}$因固定$q$消失. 剩下
$\delta S=\int\sum[(\dot q_i-\partial_{p_i}H)\delta p_i+(-\dot p_i-\partial_{q_i}H)\delta q_i]\dd t$. 各变分独立, 得$\dot q_i=\partial_{p_i}H$, $\dot p_i=-\partial_{q_i}H$. 无需把$\delta p$端点也固定; 原题least action在这里指驻值, 不保证任意路径的全局最小.
\end{solution}
\originalsection{自由相对论粒子的Legendre变换}
\begin{exercise}
原题4. $H=\sqrt{m^2c^4+c^2p^2}$, 从$L=\pp\cdot\dot{\rr}-H$求速度、能量及$L(v)$.
\end{exercise}
\begin{solution}
$\vct v=\partial_{\pp}H=c^2\pp/H$, 联立$H^2=m^2c^4+c^2p^2$给$H=\gamma mc^2$, $\pp=\gamma m\vct v$, $\gamma=(1-v^2/c^2)^{-1/2}$. 所以$L=\gamma mv^2-\gamma mc^2=-mc^2\sqrt{1-v^2/c^2}$. 小速展开为$-mc^2+mv^2/2+\cdots$, 与非相对论L只差常静能. $H$是总能量, 动能为$H-mc^2$.
\end{solution}
\originalsection{Kepler轨道及闭合条件}
\begin{exercise}
原题5. 电荷$q$被静止$-q$吸引, 平面内$L=m(\dot\rho^2+\rho^2\dot\phi^2)/2+\kappa/\rho$, $\kappa=q^2/(4\pi\eps)$. 求守恒量和径向一阶式; (a) 圆轨道周期; (b) 有界非碰撞轨道的椭圆与频率关系.
\end{exercise}
\begin{solution}
$\phi$未显含给$M=m\rho^2\dot\phi$守恒, $L$无显时给$E=m\dot\rho^2/2+M^2/(2m\rho^2)-\kappa/\rho$守恒. 所以$\dot\rho^2=2(E+\kappa/\rho)/m-M^2/(m^2\rho^2)$, 径向运动只需初值和正负支.

(a) 圆轨道$m\rho\dot\phi^2=\kappa/\rho^2$, $M^2=m\kappa\rho$, 得$T^2=4\pi^2m\rho^3/\kappa$. (b) $u=1/\rho$, 由$\dot\phi=M/(m\rho^2)$得$u''+u=m\kappa/M^2$, 撇号对$\phi$. 故$u=(m\kappa/M^2)[1+e\cos(\phi-\phi_0)]$, $e^2=1+2EM^2/(m\kappa^2)$. $E<0,M\ne0$给$e<1$椭圆, 吸引中心是焦点; 一次径向循环对应$\Delta\phi=2\pi$, 径向周期等于绕行周期, 轨道闭合. $E=0$抛物线、$E>0$双曲线不闭合; 原卷generic trajectory的椭圆陈述须加有界条件. 角速度沿椭圆并非匀速, 相等的是周期频率.
\end{solution}
